\section{Six data verbs}
\subsection{Predict, run, check, explain}

\begin{frame}{Six verbs, six decisions}
  \footnotesize
  \renewcommand{\arraystretch}{1.18}
  \begin{tabularx}{\textwidth}{@{}>{\bfseries}p{2.2cm}X@{}}
    \toprule
    Verb & Question before we run it \\
    \midrule
    \texttt{select()} & Which variables and identifiers remain? \\
    \texttt{filter()} & Which observations leave the population? \\
    \texttt{mutate()} & What new value is created, and in which unit? \\
    \texttt{group\_by()} & Which observations belong together? \\
    \texttt{summarise()} & What does one output row represent? \\
    \texttt{left\_join()} & Which keys connect the tables? \\
    \bottomrule
  \end{tabularx}
\end{frame}

\begin{frame}{Use one recurring move}
  \begin{center}
  \begin{tikzpicture}[
    node distance=0.42cm,
    processbox/.style={draw=palegray,rounded corners=2pt,text width=2.7cm,
      minimum height=1.35cm,align=center,font=\small},
    arr/.style={-{Latex[length=2mm]},color=crimson,thick}
  ]
    \node[processbox] (predict) {\textbf{Predict}\\What should change?};
    \node[processbox,right=of predict] (run) {\textbf{Run}\\Execute one small step};
    \node[processbox,right=of run] (check) {\textbf{Check}\\Did that change occur?};
    \node[processbox,right=of check] (explain) {\textbf{Explain}\\Why should it remain?};
    \draw[arr] (predict)--(run);
    \draw[arr] (run)--(check);
    \draw[arr] (check)--(explain);
  \end{tikzpicture}
  \end{center}

  \vfill
  Predict rows, columns, values, order, or unit before running the step.
\end{frame}

\begin{frame}[fragile]{\texttt{select()} defines the scope}
\begin{lstlisting}[language=R]
health <- wdi |>
  select(
    iso3c,
    country,
    year,
    income_level_current,
    gdp_per_capita_ppp,
    under5_mortality
  )
\end{lstlisting}

  \textbf{Predict:} Which columns will leave? Which identifiers need to stay?
\end{frame}

\begin{frame}[fragile]{Check after \texttt{select()}}
\begin{lstlisting}[language=R]
names(health)
glimpse(health)

stopifnot(all(
  c("iso3c", "year") %in% names(health)
))
\end{lstlisting}

  Run the key check after confirming that the contract identifiers are still
  present.
\end{frame}

\begin{frame}[fragile]{\texttt{filter()} changes the population}
\begin{lstlisting}[language=R]
health <- health |>
  filter(year >= 2000)
\end{lstlisting}

  Before running:
  \begin{itemize}
    \item How many years should remain?
    \item Could a country disappear completely?
    \item Are we removing missing indicators on purpose or by accident?
  \end{itemize}
\end{frame}

\begin{frame}[fragile]{Check after \texttt{filter()}}
\begin{lstlisting}[language=R]
health |>
  summarise(
    rows = n(),
    countries = n_distinct(iso3c),
    first_year = min(year),
    last_year = max(year)
  )
\end{lstlisting}

  Compare country and year coverage to identify which observations left the
  sample.
\end{frame}

\begin{frame}[fragile]{\texttt{mutate()} makes a claim about meaning}
\begin{lstlisting}[language=R]
health <- health |>
  mutate(
    log_income = log(gdp_per_capita_ppp)
  )
\end{lstlisting}

  The logarithm compresses the upper tail. It also changes how we interpret a
  later coefficient. The new column carries a substantive decision about scale.
\end{frame}

\begin{frame}[fragile]{Check after \texttt{mutate()}}
\begin{lstlisting}[language=R]
health |>
  summarise(
    missing_income = sum(is.na(gdp_per_capita_ppp)),
    missing_log = sum(is.na(log_income)),
    min_income = min(gdp_per_capita_ppp, na.rm = TRUE),
    max_income = max(gdp_per_capita_ppp, na.rm = TRUE)
  )
\end{lstlisting}

  What would happen if income were zero or negative?
\end{frame}

\begin{frame}[fragile]{Grouping changes the unit}
\begin{lstlisting}[language=R]
regional_year <- health |>
  group_by(region, year) |>
  summarise(
    mean_mortality = mean(under5_mortality, na.rm = TRUE),
    countries = n_distinct(iso3c),
    .groups = "drop"
  )
\end{lstlisting}

  Grouping produces one region--year row in this summary.
\end{frame}

\begin{frame}{Predict the output}
  Suppose the table has seven regions and 23 years.

  \begin{enumerate}
    \item What is the largest possible row count after the summary?
    \item Why might the actual count be smaller?
    \item Does \texttt{na.rm = TRUE} solve missing coverage?
  \end{enumerate}

  \vfill
  Take three minutes. Write the answers before we run the code.
\end{frame}
